\documentclass[10pt, letterpaper]{article}

% Packages:
\usepackage[
    ignoreheadfoot, % set margins without considering header and footer
    top=2 cm, % seperation between body and page edge from the top
    bottom=2 cm, % seperation between body and page edge from the bottom
    left=2 cm, % seperation between body and page edge from the left
    right=2 cm, % seperation between body and page edge from the right
    footskip=1.0 cm, % seperation between body and footer
    % showframe % for debugging 
]{geometry} % for adjusting page geometry
\usepackage[explicit]{titlesec} % for customizing section titles
\usepackage{tabularx} % for making tables with fixed width columns
\usepackage{array} % tabularx requires this
\usepackage[dvipsnames]{xcolor} % for coloring text
\definecolor{primaryColor}{RGB}{0, 79, 144} % define primary color
\usepackage{enumitem} % for customizing lists
\usepackage{fontawesome5} % for using icons
\usepackage{amsmath} % for math
\usepackage[
    pdftitle={David Ramírez CV},
    pdfauthor={David Ramírez Betancourth},
    pdfcreator={LaTeX with RenderCV},
    colorlinks=true,
    urlcolor=primaryColor
]{hyperref} % for links, metadata and bookmarks
\usepackage[pscoord]{eso-pic} % for floating text on the page
\usepackage{calc} % for calculating lengths
\usepackage{bookmark} % for bookmarks
\usepackage{lastpage} % for getting the total number of pages
\usepackage{changepage} % for one column entries (adjustwidth environment)
\usepackage{paracol} % for two and three column entries
\usepackage{ifthen} % for conditional statements
\usepackage{needspace} % for avoiding page brake right after the section title
\usepackage{iftex} % check if engine is pdflatex, xetex or luatex

% Ensure that generate pdf is machine readable/ATS parsable:
\ifPDFTeX
    \input{glyphtounicode}
    \pdfgentounicode=1
    \usepackage[T1]{fontenc}
    \usepackage[utf8]{inputenc}
    \usepackage{lmodern}
\fi

\usepackage[default, type1]{sourcesanspro} 

% Some settings:
\AtBeginEnvironment{adjustwidth}{\partopsep0pt} % remove space before adjustwidth environment
\pagestyle{empty} % no header or footer
\setcounter{secnumdepth}{0} % no section numbering
\setlength{\parindent}{0pt} % no indentation
\setlength{\topskip}{0pt} % no top skip
\setlength{\columnsep}{0.15cm} % set column seperation
\makeatletter
\let\ps@customFooterStyle\ps@plain % Copy the plain style to customFooterStyle
\patchcmd{\ps@customFooterStyle}{\thepage}{
    \color{gray}\textit{\small David Ramírez Betancourth - Page \thepage{} of \pageref*{LastPage}}
}{}{} % replace number by desired string
\makeatother
\pagestyle{customFooterStyle}

\titleformat{\section}{
    % avoid page braking right after the section title
    \needspace{4\baselineskip}
    % make the font size of the section title large and color it with the primary color
    \Large\color{primaryColor}
}{
}{
}{
    % print bold title, give 0.15 cm space and draw a line of 0.8 pt thickness
    % from the end of the title to the end of the body
    \textbf{#1}\hspace{0.15cm}\titlerule[0.8pt]\hspace{-0.1cm}
}[] % section title formatting

\titlespacing{\section}{
    % left space:
    -1pt
}{
    % top space:
    0.3 cm
}{
    % bottom space:
    0.2 cm
} % section title spacing

% \renewcommand\labelitemi{$\vcenter{\hbox{\small$\bullet$}}$} % custom bullet points
\newenvironment{highlights}{
    \begin{itemize}[
        topsep=0.10 cm,
        parsep=0.10 cm,
        partopsep=0pt,
        itemsep=0pt,
        leftmargin=0.4 cm + 10pt
    ]
}{
    \end{itemize}
} % new environment for highlights

\newenvironment{highlightsforbulletentries}{
    \begin{itemize}[
        topsep=0.10 cm,
        parsep=0.10 cm,
        partopsep=0pt,
        itemsep=0pt,
        leftmargin=10pt
    ]
}{
    \end{itemize}
} % new environment for highlights for bullet entries


\newenvironment{onecolentry}{
    \begin{adjustwidth}{
        0.2 cm + 0.00001 cm
    }{
        0.2 cm + 0.00001 cm
    }
}{
    \end{adjustwidth}
} % new environment for one column entries

\newenvironment{twocolentry}[2][]{
    \onecolentry
    \def\secondColumn{#2}
    \setcolumnwidth{\fill, 4.5 cm}
    \begin{paracol}{2}
}{
    \switchcolumn \raggedleft \secondColumn
    \end{paracol}
    \endonecolentry
} % new environment for two column entries

\newenvironment{threecolentry}[3][]{
    \onecolentry
    \def\thirdColumn{#3}
    \setcolumnwidth{1 cm, \fill, 4.5 cm}
    \begin{paracol}{3}
    {\raggedright #2} \switchcolumn
}{
    \switchcolumn \raggedleft \thirdColumn
    \end{paracol}
    \endonecolentry
} % new environment for three column entries

\newenvironment{header}{
    \setlength{\topsep}{0pt}\par\kern\topsep\centering\color{primaryColor}\linespread{1.5}
}{
    \par\kern\topsep
} % new environment for the header

\newcommand{\placelastupdatedtext}{% \placetextbox{<horizontal pos>}{<vertical pos>}{<stuff>}
  \AddToShipoutPictureFG*{% Add <stuff> to current page foreground
    \put(
        \LenToUnit{\paperwidth-2 cm-0.2 cm+0.05cm},
        \LenToUnit{\paperheight-1.0 cm}
    ){\vtop{{\null}\makebox[0pt][c]{
        \small\color{gray}\textit{Last updated in July 2026}\hspace{\widthof{Last updated in September 2024}}
    }}}%
  }%
}%

% save the original href command in a new command:
\let\hrefWithoutArrow\href

% new command for external links:
\renewcommand{\href}[2]{\hrefWithoutArrow{#1}{\ifthenelse{\equal{#2}{}}{ }{#2 }\raisebox{.15ex}{\footnotesize \faExternalLink*}}}


\begin{document}
    \newcommand{\AND}{\unskip
        \cleaders\copy\ANDbox\hskip\wd\ANDbox
        \ignorespaces
    }
    \newsavebox\ANDbox
    \sbox\ANDbox{}

    \placelastupdatedtext
    \begin{header}
        \fontsize{30 pt}{30 pt}
        \textbf{David Ramírez Betancourth}

        \vspace{0.3 cm}

        \normalsize
        \mbox{{\footnotesize\faMapMarker*}\hspace*{0.13cm}Manizales, Colombia}%
        \kern 0.25 cm%
        \AND%
        \kern 0.25 cm%
        \mbox{\hrefWithoutArrow{mailto:dramirezbe@unal.edu.co}{{\footnotesize\faEnvelope[regular]}\hspace*{0.13cm}dramirezbe@unal.edu.co}}%
        \kern 0.25 cm%
        \AND%
        \kern 0.25 cm%
        \mbox{\hrefWithoutArrow{tel:+57-310-680-13-31}{{\footnotesize\faPhone*}\hspace*{0.13cm}+57 3106801331}}%
        \kern 0.25 cm%
        \AND%
        \kern 0.25 cm%
        \mbox{\hrefWithoutArrow{https://linkedin.com/in/david-ramirez-betancourth-b5ba80279}{{\footnotesize\faLinkedinIn}\hspace*{0.13cm}david-ramirez-betancourth}}%
        \kern 0.25 cm%
        \AND%
        \kern 0.25 cm%
        \mbox{\hrefWithoutArrow{https://github.com/dramirezbe}{{\footnotesize\faGithub}\hspace*{0.13cm}dramirezbe}}%
    \end{header}

    \vspace{0.3 cm - 0.3 cm}

    \section{Perfil Profesional}

        \begin{onecolentry}
            Ingeniero Electrónico especializado en sistemas de radiofrecuencia (RF) y Procesamiento Digital de Señales (DSP). Experiencia en integración hardware--software desde protocolos de bajo nivel (TCP, UDP, ZMQ) y Radio Definida por Software (SDR) hasta arquitecturas cloud con sensores IoT. Desarrollo de firmware en C para adquisición y procesamiento en tiempo real, con calibración autónoma y orquestación Python. Complemento mi flujo de desarrollo con agentes de IA para automatización TDD/SDD, enfocándome en la optimización de arquitecturas embebidas de alto rendimiento.
        \end{onecolentry}

        \vspace{0.2 cm}

    \section{Educación}

        \begin{threecolentry}{\textbf{MSc}}{
            Ene 2026 - Presente
        }
            \textbf{Universidad Nacional de Colombia}, Admisión anticipada a Maestría en Automatización Industrial
            \begin{highlights}
                \item \textbf{Áreas de Investigación:} Inteligencia Artificial, Aprendizaje de Máquina aplicado a RF, Estructuras Computacionales y Optimización.
            \end{highlights}
        \end{threecolentry}

        \begin{threecolentry}{\textbf{BS}}{
            Jul 2022 - Jun 2026
        }
            \textbf{Universidad Nacional de Colombia}, Ingeniería Electrónica
            \begin{highlights}
                \item \textbf{Promedio Acumulado:} 4.2/5.0
                \item \textbf{Áreas de Enfoque:} Procesamiento Digital de Señales (DSP), Radio Definida por Software (SDR), Diseño de Hardware y Sistemas IoT.
            \end{highlights}
        \end{threecolentry}


    \section{Experiencia}

        \begin{twocolentry}{
            Manizales, Colombia

        Oct 2024 – Actualidad
        }
            \textbf{Grupo de Investigación GCPDS}, Investigador y Desarrollador
            \begin{highlights}
                \item Desarrollo de firmware C/Python para sensor IoT de monitoreo espectral basado en HackRF SDR sobre Raspberry Pi 5, con arquitectura híbrida (C para adquisición IQ en tiempo real, Python para orquestación), calibración autónoma de PPM y streaming WebRTC.
                \item Implementación de plataforma backend con API REST (FastAPI), base de datos TimescaleDB para series temporales y dashboard de visualización espectral en tiempo real.
                \item Desarrollo de aplicación móvil para despliegue y validación de modelos de Segmentación de Imágenes.
            \end{highlights}
        \end{twocolentry}

        \vspace{0.2 cm}

        \begin{twocolentry}{
            Manizales, Colombia
            
            2024
        }
            \textbf{Universidad Nacional de Colombia}, Becario de Pregrado
            \begin{highlights}
                \item Monitor académico de Señales y Sistemas y Teoría de Aprendizaje de Máquina, asesorando a grupos de 30+ estudiantes en componentes teórico-prácticos.
                \item Diseño e instrucción de talleres prácticos en Python para análisis y procesamiento de señales, integrando herramientas como NumPy, SciPy y Matplotlib.
            \end{highlights}
        \end{twocolentry}


    \section{Proyectos}
        
        \begin{twocolentry}{
            \href{https://github.com/dramirezbe/SDR-SpectrumMonitoring-Sensor}{Repo: SDR-Sensor}
        }
            \textbf{Sensor IoT para Monitoreo del Espectro}
            \begin{highlights}
                \item Desarrollo de firmware para un sistema de monitoreo espectral integrando Raspberry Pi 5 y HackRF SDR. Incluye implementación de protocolos de calibración, sincronización NTP precisa y controladores optimizados para el hardware SDR.
                \item \textbf{Herramientas:} C, Python, HackRF, Raspberry Pi.
            \end{highlights}
        \end{twocolentry}


        \vspace{0.2 cm}

        \begin{twocolentry}{
            \href{https://github.com/dramirezbe/RPI-DashWake.git}{Repo: RPI-DashWake}
        }
            \textbf{RPI-DashWake (Smart Clock)}
            \begin{highlights}
                \item Diseño de un reloj inteligente con monitoreo ambiental integrado (temperatura, humedad y calidad del aire). Cuenta con un servidor web para visualización de métricas en tiempo real y gestión programable de alarmas.
                \item \textbf{Tecnologías:} C, Python, Arduino, JavaScript.
            \end{highlights}
        \end{twocolentry}

        \vspace{0.2 cm}

        \begin{twocolentry}{
            \href{https://github.com/dramirezbe/SelfHostedContextCrafter}{Repo: ContextCrafter}
        }
            \textbf{SelfHostedContextCrafter (Orquestador de Agentes IA)}
            \begin{highlights}
                \item Equipo arquitectónico automatizado que orquesta agentes expertos basados en RAG, ejecutándose completamente en hardware local. Digiere datasheets de hardware, documentación de software y protocolos de red para generar contexto estructurado y automatizar flujos de desarrollo.
                \item \textbf{Tecnologías:} Python, RAG, LLMs locales, LaTeX.
            \end{highlights}
        \end{twocolentry}

    \section{Logros Técnicos}

        \begin{onecolentry}
            \begin{highlightsforbulletentries}
                \item \textbf{Arquitectura SDR agnóstica:} Diseño de arquitectura para una red de sensores cognitivos, integrando WebRTC para streaming en tiempo real, TimescaleDB para datos históricos y visualización de alto rendimiento en WebGL.
                \item \textbf{DSP y Calibración SDR:} Desarrollo de algoritmos de calibración iterativa (descenso de gradiente) para corregir el offset de reloj (PPM) en hardware SDR, e implementación de motores DSP personalizados en C99 usando Polyphase Filter Banks (PFB) y \textit{liquid-dsp}.
                \item \textbf{Telemetría y Profiling en Pi 5:} Ejecución de pruebas de estrés de hardware y perfilado de energía (configuraciones PD/PPS) en Raspberry Pi 5, identificando y mitigando caídas de voltaje bajo cargas concurrentes de LTE y procesamiento de radiofrecuencia.
                \item \textbf{Automatización con IA (TDD/SDD):} Automatización de flujos de desarrollo mediante agentes de IA con metodologías TDD y SDD, generando código, tests y documentación de forma predictiva. Reducción significativa de tareas repetitivas de boilerplate y validación.
                \item \textbf{Orquestador de Agentes (SelfHostedContextCrafter):} Desarrollo de un sistema que orquesta agentes expertos basados en RAG en hardware local, capaz de digerir datasheets de hardware y documentación técnica para generar contexto estructurado y automatizar tareas de revisión y despliegue de código.
            \end{highlightsforbulletentries}
        \end{onecolentry}

        \vspace{0.2 cm}
    
    \section{Tecnologías}

        \begin{onecolentry}
            \textbf{Lenguajes:} C99, C++, Python, JavaScript, TypeScript, Verilog, SQL
        \end{onecolentry}

        \vspace{0.2 cm}

        \begin{onecolentry}
            \textbf{Frameworks y Plataformas:} FastAPI, Django, React, Vite, Node.js, GNU Radio, SoapySDR, HackRF One, FPGAs (Gowin EDA)
        \end{onecolentry}

        \vspace{0.2 cm}

        \begin{onecolentry}
            \textbf{Infraestructura y Herramientas:} Linux (systemd, Bash), Docker, CMake, ZMQ, WebRTC, TimescaleDB, Git, Sphinx, Doxygen, NTP, Raspberry Pi, GPS/LTE
        \end{onecolentry}

    \section{Idiomas}

        \begin{onecolentry}
            \textbf{Español:} Nativo \quad|\quad \textbf{Inglés:} Intermedio (B1)
        \end{onecolentry}


\end{document}